% ============================================================
%  文档类选项（在 main.tex 的 \documentclass[...]{tongjithesis} 中设置）
%  field=science    理工科（工科类、理科类，默认）— 章节编号：1 / 1.1 / 1.1.1
%  field=humanities 文科  （人文、法学、外语、艺术）        — 章节编号：一、/（一）/ 1.
%  degree=bachelor  学士学位（默认）
%  degree=master   硕士学位
%  degree=doctor   博士学位
% ============================================================

% ============================================================
%  封面信息
% ============================================================
\school{国豪书院}
\major{计算机科学与技术}
\student{2251251}{姚雨萱}
\thesistitle{基于多模态大语言模型的胶质母细胞瘤候选药物分子生成与优化研究}{}
\thesistitleeng{Glioblastoma Drug Candidate Molecule Generation and Optimization Based on Multi-modal Large Language Model}{}
\thesisadvisor{关佶红\quad 教授}
\thesisdate{2026}{5}{10} % 使用当前日期; 也可以自定义日期

% ============================================================
%  信息说明页
% ============================================================

% 成果类型（三选一）：
%   thesis      — 毕业论文，摘要
%   design      — 毕业设计（软件/创意/建筑等非工程设计类），摘要
%   engineering — 毕业设计（工程设计类），设计总说明
\infotype{thesis}

% 内容简述（请用300字以内简要概述）
\infoabstract{%
本文针对胶质母细胞瘤（Glioblastoma, GBM）新药早期发现中候选分子生成难、多目标约束平衡复杂等问题，提出了一种名为 GBM-PharmaMind 的候选药物分子生成与优化框架。该框架以多模态大语言模型为基础，通过低秩适配（Low-Rank Adaptation, LoRA）微调进行胶质母细胞瘤领域知识适配，结合图扩散变换器实现分子图生成与编码，并引入思维链（Chain-of-Thought, CoT）与思维树（Tree-of-Thoughts, ToT）协同推理，引导模型在靶点信息、血脑屏障穿透性和成药性约束下生成候选分子。生成结果结合分子对接、成药性预测和结构风险检测进行评估与反馈，实现迭代优化。在代表性靶点上的实验结果表明，该框架能够生成具有一定靶点匹配潜力、血脑屏障穿透性和成药性特征的候选分子；对比实验表明，该框架在高分候选分子质量方面具有优势；消融实验说明领域知识适配、协同推理和评估反馈等模块均有助于提升生成质量。本文研究可为胶质母细胞瘤创新药物的早期虚拟筛选提供计算辅助思路。
}

% 毕业设计：图纸数量与正文字数（成果类型为 design 或 engineering 时填写，两者须同时填写）
% \infodrawings{5}
% \infowordcount{15000}

% 毕业论文：正文总字数（成果类型为 thesis 时填写）
\infothesiswords{27120}

% 随附资料（每条用 \item 开头；不填则显示空白行）
\infomaterials{
  \item GBM-PharmaMind框架源代码及简要使用说明
}
